\section{Summary of Achievements}

This project successfully demonstrates the integration of multiple sensor modalities (RGB camera and depth camera) for pedestrian detection and tracking in an autonomous driving scenario. The key accomplishments include:

\begin{enumerate}
    \item \textbf{Perception Pipeline}: Implemented a complete perception system that fuses RGB and depth sensor data for robust object detection and tracking.
    
    \item \textbf{Object Detection}: Integrated YOLOv8 for real-time object detection with persistent tracking IDs across frames.
    
    \textbf{Object Tracking}: Developed a custom 2D Kalman filter for tracking pedestrian trajectories and estimating velocities relative to the ego vehicle.
    
    \item \textbf{Collision Assessment}: Implemented a TTC (Time-to-Collision) based collision warning system that can distinguish between:
    \begin{itemize}
        \item Safe crossing (pedestrian not in trajectory)
        \item Warning phase (pedestrian entering trajectory, possible collision)
    \end{itemize}
    
    \item \textbf{Autonomous Control}: Created a velocity-based control system that can brake automatically when collision is imminent.
    
    \item \textbf{Scenario Support}: Provided a flexible scenario configuration system supporting multiple test cases (colliding vs. non-colliding pedestrians).
    
    \item \textbf{Real-time Visualization}: Implemented a real-time visualization tool showing detections, predictions, and trajectories.
\end{enumerate}

\section{Key Design Choices}

The project architecture was designed around the principle of separation of concerns:

\begin{enumerate}
    \item \textbf{Core vs. Utils}: Core modules contain the main logic, while utility modules provide reusable helper functions.
    
    \item \textbf{Simulation Loop}: The main loop clearly separates camera callback processing (asynchronous) from the main simulation loop (synchronous).
    
    \item \textbf{Configuration}: YAML-based configuration allows easy modification of simulation parameters without code changes.
\end{enumerate}

\section{Challenges Encountered}

During development, we encountered several challenges:

\begin{enumerate}
    \item \textbf{CARLA Performance}: Setting up CARLA with multiple sensors can be computationally expensive. We optimized this by using a lightweight YOLOv8 model and limiting simulation FPS.
    
    \item \textbf{Depth Calibration}: Achieving accurate depth measurements required careful sensor calibration. Small misalignments can lead to significant positioning errors.
    
    \item \textbf{YOLOv8 Latency}: The YOLOv8 nano model achieved good speed ($\sim~$40-60 FPS on our test GPU) but still limited the overall simulation to $\sim~20$ FPS due to other processing.
    
    \item \textbf{Pedestrian Model}: CARLA's \texttt{walker} actor doesn't have the same mesh as real vehicles, which can cause minor visual inconsistencies but doesn't affect functionality.
    
    \item \textbf{Ego Vehicle Control}: The simple velocity controller doesn't handle complex scenarios (lane changes, obstacle avoidance). A full autonomous driving stack would be significantly more complex.
\end{enumerate}

\section{Results Summary}

The system successfully:

\begin{enumerate}
    \item Detected pedestrians with YOLOv8 (model: yolov8n.pt)
    \item Tracked objects across frames using a Kalman filter
    \item Estimated trajectories and future positions
    \item Calculated TTC for collision assessment
    \item Triggered braking when collision risk is detected
    
    \item Correctly classified collision scenarios:
    \begin{itemize}
        \item No collision (pedestrian far from path)
        \item Warning (pedestrian in trajectory, $<$5s until impact)
    \end{itemize}
\end{enumerate}

\section{Performance Metrics}

\begin{table}[h]
    \centering
    \begin{tabular}{|l|l|l|l|}
        \hline
        Component & Metric & Value & Notes \\
        \hline
        YOLOv8 inference & FPS & 45 & On AMD 7800 XT, yolov8n.pt \\
        \hline
        Kalman update & Time per frame & $\sim$1-2ms & Negligible \\
        \hline
        TTC calculation & Time per call & $<$100µs & Extremely fast \\
        \hline
        Total system & Max FPS & 20 & Configurable via simulation.fps \\
        \hline
        Memory usage & Peak RAM & $~$2-3GB & Dominated by CARLA \\
        \hline
        Simulation window & Latency & $\sim$50-100ms & Acceptable for real-time \\
        \hline
    \end{tabular}
    \caption{Performance metrics for key components}
\end{table}

\section{Future Improvements}

Based on our experience, the following improvements could enhance the system:

\begin{enumerate}
    \item \textbf{Better Pedestrian Model}: Replace CARLA walkers with a more realistic pedestrian model (e.g., from CARLA worldgen) that includes:
    \begin{itemize}
        \item Looking behavior (pedestrian looking at car before crossing)
        \item Uncertainty in decision making
        \item Multiple pedestrian types (elderly, children, etc.)
    \end{itemize}
    
    \item \textbf{Improved Vehicle Control}: Replace simple velocity control with a full autonomous driving stack:
    \begin{itemize}
        \item Lane keeping
        \item Obstacle avoidance
        \item Adaptive cruise control
        \item Traffic light compliance
    \end{itemize}
    
    \item \textbf{Enhanced Sensor Fusion}: Integrate LiDAR and radar data for more robust perception:
    \begin{itemize}
        \item LiDAR for accurate 3D positioning
        \item Radar for velocity estimation
        \item Camera for semantic understanding
    \end{itemize}
    
    \item \textbf{Advanced Tracking}: Implement a more sophisticated tracking system:
    \begin{itemize}
        \item Multi-object tracking with ID consistency
        \item Track spawning and deletion
        \item Re-ID for occluded objects
        \item EKF or UKF for non-linear dynamics
    \end{itemize}
    
    \item \textbf{Better Visualization}: Add AR-style visualization:
    \begin{itemize}
        \item Project predictions onto 3D scene
        \item Overlay confidence levels
        \item Show multiple sensor modalities
        \item HUD display for driver warnings
    \end{itemize}
    
    \item \textbf{Multi-Agent Simulation}: Simulate multiple vehicles with realistic traffic:
    \begin{itemize}
        \item Multiple autonomous vehicles
        \item Realistic human-driven vehicles
        \item Traffic rules compliance
    \end{itemize}
    
    \item \textbf{Simulation Scenarios}: Add more diverse scenarios:
    \begin{itemize}
        \item Different road layouts (intersection, highway, etc.)
        \item Varying weather conditions
        \item Day/night cycles
        \item Different pedestrian behaviors (crossing, jaywalking, etc.)
    \end{itemize}
    
    \item \textbf{Testing Framework}: Develop a comprehensive test framework:
    \begin{itemize}
        \item Unit tests for all modules
        \item Integration tests for end-to-end scenarios
        \item Regression tests
        \item Performance benchmarks
    \end{itemize}
\end{enumerate}

\section{Concluding Remarks}

This project demonstrates the feasibility of a sensor-fused perception system for autonomous driving, specifically for pedestrian crossing scenarios. By combining YOLOv8 object detection with a Kalman-based tracking system and TTC-based collision assessment, we achieved a functional prototype that can detect and warn about potential collisions with pedestrians.

The system architecture is modular and extensible, making it suitable for:
\begin{itemize}
    \item Further development of autonomous driving capabilities
    \item Integration with other perception systems
    \item Use as a teaching tool for autonomous driving concepts
    \item Benchmarking of different perception algorithms
\end{itemize}
